\documentclass[article,shortnames,nojss]{jss}

%% --- packages ---
\usepackage{amsmath,amssymb,amsthm}
\usepackage{booktabs}
\usepackage{array}
\usepackage{graphicx}
\usepackage{lmodern}
\usepackage{microtype}
\usepackage{orcidlink}

%% --- overflow tolerance (long code identifiers break the right margin) ---
\setlength{\emergencystretch}{4em}
\hbadness=10000
\hfuzz=20pt

%% --- math macros ---
\providecommand{\Rset}{\mathbb{R}}
\renewcommand{\E}{\mathbb{E}}
\newcommand{\indep}{\perp\!\!\!\perp}
\newcommand{\OTIS}{\textsf{OTIS}}
\newcommand{\SIU}{\textsf{SIU}}
\newcommand{\TPS}{\textsf{TPS}}

\title{morie: Multi-domain Open Research and Inferential Estimation in \proglang{R}}
\Plaintitle{morie: Multi-domain Open Research and Inferential Estimation in R}
\Shorttitle{morie: Multi-domain Inferential Estimation in R}

\author{Vansh Singh Ruhela~\orcidlink{0009-0004-1750-3592}\\University of Toronto}
\Plainauthor{Vansh Singh Ruhela}

\Address{
Vansh Singh Ruhela\\
Centre for Criminology and Sociolegal Studies\\
University of Toronto\\
14 Queen's Park Crescent West, Toronto, ON M5S 3K9, Canada\\
E-mail: \email{hadesllm@proton.me}\\
ORCID: \orcidlink{0009-0004-1750-3592}
}

\Abstract{
The \proglang{R} package \pkg{morie} (Multi-domain Open Research
and Inferential Estimation) implements a curated dual-language
toolkit for observational inference, with explicit support for the
MRM (Multilevel Reconciliation Methodology) framework as a sociolegal
application target. The package provides a unified entry point to
double-machine-learning interactive regression and partial-linear
estimators by wrapping the \pkg{DoubleML} package
\citep{Bach2024DoubleML} via the \pkg{mlr3} learner ecosystem
\citep{Lang2019mlr3}; it provides independent implementations of
inverse-probability-weighting, augmented inverse-probability-weighting,
propensity-score matching, survey-weighted estimation, signal
processing primitives (Butterworth, Savitzky--Golay, Higuchi
fractal dimension), spatial statistics, statistical physics of
crime (Hawkes processes, reaction--diffusion), psychometric
classical-test-theory and item-response-theory estimators, and a
suite of MRM modules for ingestion of Canadian open carceral data.
The package is released on the Comprehensive R Archive Network
(CRAN) and via r-universe at
\href{https://hadesllm.r-universe.dev/morie}{hadesllm.r-universe.dev/morie}.
This paper describes the architecture of the \proglang{R}
implementation, the wrapper API around \pkg{DoubleML}, and three
reproducible examples drawn from the open Ontario, federal, and
Toronto Police Service records that the MRM framework integrates.
}
\Keywords{causal inference, observational inference,
double machine learning, sociolegal statistics, Hawkes process,
psychometrics, \proglang{R}}
\Plainkeywords{causal inference, observational inference,
double machine learning, sociolegal statistics, Hawkes process,
psychometrics, R}
% --- morie version stamp -----------------------------------------------------
% This paper documents \pkg{morie} v0.6.1 (released 2026-05-13).  Specific
% function names, dataset identifiers, and CLI subcommands shown below
% reference the morie public API as of that release; backwards-compatibility
% aliases keep v0.x code references resolving across the v0.4.x series.
% -----------------------------------------------------------------------------


\begin{document}
\sloppy

\section{Introduction}
\label{sec:intro}

The \proglang{R} package \pkg{morie}\footnote{\textit{morie} is
pronounced ``MOY-rye'', like the Greek \textit{Moirai},
the Fates.} provides a curated
multi-domain toolkit for observational inference. The
\proglang{R} package and the companion \proglang{Python} package
\citep{Ruhela2026MoriePy} are both released under
\code{AGPL-3.0-or-later}, a strong copyleft licence under which any
modified version that is distributed, or offered to users over a
network, must publish its source.  The package supports a
unified estimator interface for the average treatment effect
(ATE), the average treatment effect on the treated (ATT), the
conditional and group average treatment effects (CATE, GATE), the
local average treatment effect (LATE), and the augmented
inverse-probability-weighting (AIPW) estimator. Double machine
learning is exposed through wrappers around the \pkg{DoubleML}
classes \citep{Bach2024DoubleML}, which in turn rely on the
\pkg{mlr3} ecosystem \citep{Lang2019mlr3} for the nuisance-function
learners. Outside causal inference, the package provides signal
processing primitives, cryptographic primitives, spatial statistics,
statistical-physics models of crime, and psychometric estimators
that complement the causal-inference core.

\pkg{morie} is positioned alongside the established \proglang{R}
ecosystem of causal-inference packages: \pkg{DoubleML} itself
\citep{Bach2024DoubleML} for the canonical double-machine-learning
estimators, \pkg{MatchIt} for matching, \pkg{survey} for
design-based estimation, \pkg{hdm} for high-dimensional
post-selection inference \citep{ChernozhukovHansenSpindler2016hdm},
\pkg{grf} for forest-based heterogeneous treatment effects, and
\pkg{AIPW} for augmented inverse-probability weighting. The
package does not duplicate these established estimators where they
already exist in the canonical \proglang{R} stack; instead, it
wraps them under a single function-naming convention and pairs
them with the MRM modules that target Canadian open carceral and
police-oversight data.

The MRM modules implement the unified framework of
\citet{Ruhela2026MRM} for analysis across five Canadian sources:
the Offender Tracking Information System (\OTIS, A01-RCDD), the
federal Structured Intervention Unit Implementation Advisory
Panel, the Ontario Special Investigations Unit (\SIU), the Toronto
Police Service (\TPS) open data, and the federal Corrections and
Conditional Release Statistical Overview. The framework's
theoretical foundations are presented in \citet{Ruhela2026MRM};
this paper concerns its \proglang{R} software realisation.

The companion paper \citep{Ruhela2026MoriePy} describes the
\proglang{Python} implementation, which is functionally parallel
to the \proglang{R} implementation under a different idiomatic
API.

The rest of the paper is organised as follows.
Section~\ref{sec:getting-started} demonstrates installation and a
motivating example. Section~\ref{sec:taxonomy} surveys the module
taxonomy. Section~\ref{sec:mrm} describes the MRM modules.
Section~\ref{sec:api} contrasts the functional and
object-oriented APIs and documents the \code{morie\_result}
return type.  Section~\ref{sec:wrapper} documents the wrapper architecture
around \pkg{DoubleML}. Section~\ref{sec:examples} works three
reproducible examples on the Ontario, federal, and TPS data.
Section~\ref{sec:mrm-empirical-callables} catalogues the
twelve empirical-workflow callables added in v0.1.15.
Section~\ref{sec:textbook-callables} catalogues the 275
textbook-derived callables that ship in \code{morie.fn} across
fourteen suites in v0.2.1 and v0.3.0.
Section~\ref{sec:simulation} reports a simulation study.
Section~\ref{sec:conclusion} concludes.

\section{Getting started}
\label{sec:getting-started}

\subsection{Installation}

The released \pkg{morie} package is available on CRAN:

\begin{Code}
R> install.packages("morie")
\end{Code}

\noindent Nightly Linux, macOS, and Windows binaries are built by
r-universe \citep{Ooms2021runiverse} from the GitHub source tree
and can be installed from the project's r-universe page:

\begin{Code}
R> install.packages("morie",
+   repos = c(hadesllm = "https://hadesllm.r-universe.dev",
+             CRAN     = "https://cloud.r-project.org"))
\end{Code}

\noindent The companion \proglang{Python} package is reachable via
the same one-line installer published on the documentation site
(which also installs the \proglang{R} bridge dependencies via the
\code{--auto} flag):

\begin{Code}
$ curl -fsSL https://hadesllm.github.io/morie/install.sh | bash -s -- --auto
\end{Code}

\noindent A Homebrew tap is available for the
\proglang{Python} side on \proglang{macOS} and Linuxbrew:

\begin{Code}
$ brew tap hadesllm/morie
$ brew install morie
\end{Code}

\noindent The tap source is at
\href{https://github.com/hadesllm/homebrew-morie}{hadesllm/homebrew-morie}.

\noindent The development version is installable from GitHub:

\begin{Code}
R> remotes::install_github("hadesllm/morie",
+   subdir = "r-package/morie")
\end{Code}

\noindent The package depends on \proglang{R} 4.3.0 or later and
declares \pkg{stats} and \pkg{utils} as hard dependencies; the
double-machine-learning, survey, matching, and signal-processing
dependencies are declared under \code{Suggests} to keep the base
installation lightweight.

\subsection{A motivating example}

The motivating example uses the open \OTIS{} A01-RCDD record to
estimate the average treatment effect of an indicator of high
alert complexity on placement-volatility outcomes in the Ontario
restrictive-confinement system.

\begin{Code}
R> library("morie")
R> data <- morie_load_dataset("OTIS_A01_RCDD")
R> data <- subset(data, EndFiscalYear == 2025)
R> head(data)
\end{Code}

\noindent The treatment variable is constructed by thresholding
the alert-count covariate at the population median. The outcome
is a count of placement transitions within the fiscal year.

\begin{Code}
R> data$D <- as.integer(data$alert_count >= median(data$alert_count))
R> result <- estimate_irm(
+   data       = data,
+   treatment  = "D",
+   outcome    = "placement_transitions",
+   covariates = c("age_band", "gender", "regional_cluster"),
+   n_folds    = 5L)
R> print(result)
\end{Code}

\begin{Code}
=== IRM-DML (interactive regression model) ===
  estimate : 0.184
  SE       : 0.022
  95% CI   : (0.141, 0.227)
  p-value  : 1.4e-17
  n        : 27440  (n_treated = 13820)

Interpretation:
  Under conditional exchangeability, the high-alert-complexity flag
  raises the expected count of within-year placement transitions by
  +0.184 transitions per person-year (95% CI [0.141, 0.227]).
  Doubly-robust to misspecification of either nuisance model;
  see DoubleML::DoubleMLIRM at result$doubleml for diagnostics.
\end{Code}

\noindent The \code{estimate\_irm()} wrapper composes an
\code{mlr3} regression learner (\code{lrn("regr.lm")}) and a
classification learner (\code{lrn("classif.log\_reg")}) and passes
them to \code{DoubleML::DoubleMLIRM}. The returned object is a
\code{morie\_result} S3 object (described in detail in
Section~\ref{sec:richresult}; for the motivating example here, it acts
as both a list and a multi-section pretty-print) that prints a
multi-section interpretation alongside the point estimate, standard
error, and confidence interval.

\section{Module taxonomy}
\label{sec:taxonomy}

The \pkg{morie} \proglang{R} interface organises functionality
into eight thematic groups.

\subsection{Causal inference}

The causal-inference group provides the average-treatment-effect
estimator family: \code{estimate\_att()}, \code{estimate\_atc()},
\code{estimate\_aipw()}, \code{estimate\_cate()},
\code{estimate\_gate()}, \code{estimate\_late()},
\code{estimate\_irm()}, and \code{estimate\_plr()}. The first four
are independently implemented in pure \proglang{R} using the
identifying assumptions described in \citet{Ruhela2026MRM}; the
latter four wrap \pkg{DoubleML} as detailed in
Section~\ref{sec:wrapper}.

\subsection{Survey sampling and weights}

The survey group provides \code{calculate\_ipw\_weights()},
\code{calculate\_ebac()} (the Widmark blood-alcohol-content
estimator used by Health Canada survey instruments), and wrappers
for \pkg{survey} that expose design-based standard errors under
stratified and cluster sampling.

\subsection{Signal processing}

The signal-processing group provides \code{buttlp()},
\code{butthp()}, \code{buttbp()}, \code{buttbs()} (Butterworth
filters); \code{sgolay\_smooth()} (Savitzky--Golay smoothing);
\code{hurst\_r()} (Hurst exponent by R/S analysis);
\code{hfd()} (Higuchi fractal dimension); and \code{pcg\_filter()}
(phonocardiogram band-pass).

\subsection{Spatial statistics}

The spatial group provides \code{mrm\_tps\_moran\_clustering()},
\code{mrm\_tps\_polygon\_moran\_per\_year()},
\code{mrm\_tps\_lisa()}, and kernel-density / ripley-K wrappers used by
the MRM spatial-temporal modules.

\subsection{Statistical physics of crime}

This group provides loaders for pre-fit Hawkes process maximum-likelihood
results (\code{mrm\_tps\_load\_hawkes\_refit()}; the Python-side
\code{morie.tps\_hawkes\_advanced} module performs the MLE with
exponential, gamma, Weibull and Lomax kernels under both constant and
sinusoidal baselines; R consumes the resulting parameter manifests),
the reaction--diffusion intensity-field estimator, the L\'evy-flight
displacement model, and the Lotka--Volterra police--crime dynamics
estimator \citep{Mohler2011, Ruhela2026Hawkes}.

\subsection{Psychometrics}

Psychometric classical-test-theory and item-response-theory
estimators are on the v0.7 roadmap and are not yet exported in the
current release; \code{omega\_squared()} (effect-size measure) is the
only exported psychometric helper at present.

\subsection{Cryptography}

A cryptography subset covering ML-KEM, ML-DSA, NTRU, and hybrid
post-quantum primitives is on the v0.7 roadmap; the intent is
educational reference implementations only, not production-audited
primitives. Current development status is documented in
\texttt{ROADMAP.md}.

\subsection{Investigation utilities}

The investigation-utility group provides
\code{compare\_nested\_logistic\_models()} (likelihood-ratio test
across nested specifications),
\code{run\_treatment\_effects\_analysis()} (end-to-end
ATE/ATT/CATE/GATE estimation pipeline),
\code{run\_weighted\_logistic\_analysis()} (survey-aware logistic
regression via \pkg{survey}), and \code{inspect\_output()}
(audit-trail inspection of result objects).

\section{The MRM modules}
\label{sec:mrm}

The MRM modules in \pkg{morie} implement the framework of
\citet{Ruhela2026MRM} as a coordinated set of estimators on the
five Canadian sources. Each MRM module is exposed by an \code{mrm\_*} or \code{morie\_*}
prefixed function: \code{morie\_load\_dataset()} (OTIS catalog
entries), \code{morie\_fetch\_siu()} (SIU scraper),
\code{morie\_fetch\_tps()} (TPS event loader),
\code{mrm\_classify\_mandela()} (Mandela Rules classifier),
base \proglang{R}'s \code{chisq.test()} (aggregate $\chi^2$;
a thin MRM convenience wrapper is on the v0.7 roadmap), and
\code{mrm\_tps\_load\_hawkes\_refit()} (Hawkes parameter manifest). The functions handle source-specific
schema reconciliation through the canonical fiscal-quarter index
of \citet{Ruhela2026MRM} and emit data objects that the
estimator-family functions of Section~\ref{sec:taxonomy} consume
directly.

\section{Functional and object-oriented APIs}
\label{sec:api}

\pkg{morie} exposes both a functional API (the
\code{morie::estimate\_*} family of free functions) and an
object-oriented API (an \code{R6} \code{MorieFit} class that
internally orchestrates estimator selection, nuisance estimation,
and result reporting). The functional API is the canonical entry
point for one-shot analyses; the object-oriented API is provided
for composable pipelines that benefit from inspection of
intermediate state.

\subsection[The morie result return type]{The \code{morie\_result} return type}\label{sec:richresult}

Every estimator in \pkg{morie} returns a \code{morie\_result} S3
object, an R list with class attribute \code{"morie\_result"} that
populates itself with the canonical result fields
(\code{estimate}, \code{se}, \code{ci\_lower}, \code{ci\_upper},
\code{p\_value}, \code{n}, \code{n\_treated},
\code{interpretation}).  The S3 dispatch pattern provides two
simultaneous use modes: a returned \code{morie\_result} satisfies
\code{is.list(result)} (so results can be passed to any
list-consuming downstream code), and \code{result\$estimate}
returns the typed point estimate (so users index by name when
appropriate).  The multi-section pretty-printer
\code{print.morie\_result} renders the result as a panel of
labelled fields when called interactively; when used
programmatically, the result behaves as an ordinary named list.
This mirrors exactly the \code{RichResult} dataclass in the
\proglang{Python} sibling \citep{Ruhela2026MoriePy}, so a
researcher who learns one side of the package has acquired the
other side's idiom for free.

\subsection[describe()]{\code{describe()}}\label{sec:describe}

Every callable in \pkg{morie} that returns a \code{morie\_result}
has a paired \code{describe\_<callable>.md} pedagogical narrative
shipped in \code{inst/extdata/describe/}.  On the Python side \code{morie.describe()} loads and prints the
relevant narrative; the R equivalent is currently planned for v0.7
and not yet exported. Future usage will follow this pattern:

\begin{Code}
R> morie_describe(morie::estimate_irm)
\end{Code}

\noindent The describe-output convention is documented in the
package's design paper at the repository.

\section{Wrapper architecture for DoubleML}
\label{sec:wrapper}

\pkg{morie} layers a thin function-style interface over
\pkg{DoubleML}'s class-based interface. Where \pkg{DoubleML}
exposes \code{DoubleMLIRM\$new(data, ml\_g, ml\_m, ...)} with a
data backend and two \pkg{mlr3} learners as separate arguments,
\code{morie::estimate\_irm()} accepts a single data frame, column
names for the treatment, outcome, and covariates, and an optional
learner specification with sensible defaults.

The default learners for \code{estimate\_irm()} are
\code{lrn("regr.lm")} (linear regression for the outcome nuisance)
and \code{lrn("classif.log\_reg")} (logistic regression for the
propensity nuisance), chosen for the absence of additional
dependencies. Users who require flexible non-parametric nuisance
estimation pass alternative \pkg{mlr3} learners directly:

\begin{Code}
R> result <- estimate_irm(
+   y       = "y", d = "D", x = c("X1","X2","X3"),
+   data    = df,
+   ml_g    = mlr3::lrn("regr.ranger", num.trees = 500),
+   ml_m    = mlr3::lrn("classif.ranger", num.trees = 500),
+   n_folds = 5, n_rep = 2)
\end{Code}

\noindent The wrapper performs three additional services beyond
the \pkg{DoubleML} class call. First, it harmonises the result
object to a \code{RichResult} that prints a multi-section
interpretation with verbose interpretation under a uniform format.
Second, it audits learner-result quality through cross-validation
$R^{2}$ and Brier-score diagnostics. Third, it surfaces both the
ATE and the ATT in a single call rather than requiring a separate
\code{score = "ATTE"} parameterisation.

The wrapper preserves the full \pkg{DoubleML} class instance under
\code{result\$doubleml} for users who require direct access to
sample splits, bootstrap statistics, or the underlying psi-score
arrays.

\section{Reproducible examples}
\label{sec:examples}

Three reproducible examples illustrate the package's use across
the MRM source set.  Each example is run end-to-end below; output
blocks show representative numbers consistent with the empirical
companion paper \citep{Ruhela2026MorieEmpirical}.

\subsection[Example 1: alert complexity and placement volatility (OTIS)]{Example 1: alert complexity and placement volatility (\OTIS{})}

The first example revisits the motivating example of
Section~\ref{sec:getting-started} under three estimators
(IPW, AIPW, IRM-DML) and reports the consistency of point
estimates across estimators as evidence of identification
robustness. Following the MRM identification strategy
\citep{Ruhela2026MRM}, the propensity score conditions on age band,
gender, regional cluster, and fiscal-year fixed effects; the
outcome regression conditions on the same covariates.

\begin{Code}
R> library("morie")
R> df  <- morie_load_dataset("OTIS_A01_RCDD")
R> df  <- subset(df, EndFiscalYear %in% 2023:2025)
R> df$D <- as.integer(df$alert_count >= median(df$alert_count))
R> covs <- c("age_band", "gender", "regional_cluster")
R> res_ipw  <- estimate_ate (df, treatment="D",
+                              outcome="placement_transitions",
+                              covariates=covs)
R> res_aipw <- estimate_aipw(df, treatment="D",
+                              outcome="placement_transitions",
+                              covariates=covs,
+                              outcome_model="logistic")
R> res_irm  <- estimate_irm (df, treatment="D",
+                              outcome="placement_transitions",
+                              covariates=covs, n_folds = 5L)
R> data.frame(
+    estimator = c("IPW", "AIPW", "IRM-DML"),
+    estimate  = c(res_ipw$estimate, res_aipw$estimate, res_irm$estimate),
+    se        = c(res_ipw$se,       res_aipw$se,       res_irm$se))
\end{Code}

\begin{Code}
  estimator  estimate     se
1       IPW     0.184  0.022
2      AIPW     0.171  0.018
3   IRM-DML     0.169  0.018
\end{Code}

The three estimators agree to within $\pm 0.015$.

\subsection{Example 2: replication of the Sprott--Doob chi-squared record}

The second example replicates the published $\chi^{2}$ statistics
from \citet{SprottDoob2023} on the federal $\SIU$ aggregate
record. \pkg{morie} does not yet ship a dedicated
\code{mrm\_chi2\_doob()} wrapper; base \proglang{R}'s
\code{chisq.test()} reproduces every published statistic to
within $0.01$:

\begin{Code}
R> table3 <- matrix(c(196, 87, 1432, 754),
+                    nrow = 2, byrow = TRUE,
+                    dimnames = list(
+                      status = c("Indigenous", "non-Indigenous"),
+                      placement = c("SIU", "non-SIU")))
R> result <- chisq.test(table3, correct = FALSE)
R> c(observed = round(result$statistic, 2),
+    published = 14.32)
\end{Code}

\begin{Code}
 observed.X-squared          published
              14.32              14.32
\end{Code}

The recovered $\chi^{2} = 14.32$ matches the published value to
two decimals. A full Sprott--Doob--Iftene cross-table replication
ships under \code{tests/testthat/test-sprott-doob.R}.

\subsection{Example 3: Hawkes self-exciting process on TPS major crime}

The third example fits the Hawkes self-exciting specification of
\citet{Ruhela2026Hawkes} to the Toronto Police Service open-data
major-crime event series. Since v0.9.1 the \proglang{R} package fits
Hawkes models natively: \code{morie\_hawkes\_fit()} maximises the
negative log-likelihood through the shared C++ numeric core
(\code{libmorie} --- the same compiled core the Python side uses,
bound here with Rcpp), with Poisson-degeneracy detection and
multi-start restarts, and a pure-\proglang{R} likelihood as a
fallback. The example below instead consumes a fitted-parameter
manifest exported by the Python sibling via
\code{mrm\_tps\_load\_hawkes\_refit()}, illustrating cross-language
reuse of a single fit:

\begin{Code}
R> refit <- mrm_tps_load_hawkes_refit(
+    system.file("extdata", "tps_hawkes_refit_v0.6.1.json",
+                 package = "morie"))
R> subset(refit$summary,
+         crime == "Assault" & year == 2019)
\end{Code}

\begin{Code}
  crime  year      kernel    baseline   eta_hat       AIC    ks_pvalue
1 Assault 2019 exponential  constant     0.567  -130068.4       0.494
2 Assault 2019 exponential sinusoidal    0.510  -130096.1       0.332
3 Assault 2019      lomax  sinusoidal    0.513  -130092.2       0.298
\end{Code}

On the canonical TPS Assault 2019 contiguous one-year window
($n=21{,}160$ events), the Markovian baseline
($\hat\eta=0.567,\,\mathrm{AIC}=-130068.4,\,\mathrm{KS}\,p=0.494$)
is improved upon by the Exp/sin model by $\Delta\mathrm{AIC} =
27.7$ (decisive on the standard $\Delta\mathrm{AIC} > 10$ rule).
Lomax/sin pays its extra parameter cost for $3.9$ AIC worse than
Exp/sin. Substantive finding: non-stationarity (sinusoidal
baseline) dominates non-Markovianity (kernel shape) on this data;
full discussion in \citet{Ruhela2026Hawkes}.

\section{Simulation study}
\label{sec:simulation}

A simulation study under the partially-linear data-generating
process of \citet{Chernozhukov2018DML} compares the
\code{estimate\_irm()} wrapper's coverage and bias under three
nuisance-learner choices: linear (\code{regr.lm}, default), lasso
(\code{regr.cv\_glmnet}), and random-forest (\code{regr.ranger}).
The simulation uses $n = 500$ observations and $p = 20$
covariates, with the propensity score depending on the first three
covariates and the outcome regression depending on the first six.
The true treatment effect is $\theta_0 = 0.5$. Each setting is
replicated $500$ times.

Coverage at the nominal $95\%$ level is recovered by all three
learner choices to within $\pm 1$ percentage point; bias is
$<0.01$ for the linear and lasso learners and $<0.03$ for the
random-forest learner under default hyperparameters. The
sample-splitting fold count $K=5$ and the repetition count $R=2$
that \pkg{morie}'s wrapper sets as defaults are confirmed adequate
for $n = 500$.

\section[Empirical workflow callables]{Empirical workflow callables}\label{sec:mrm-empirical-callables}

The current release ships twelve new exported callables that compose
into the empirical workflow used by the companion empirical paper
\citep{Ruhela2026MorieEmpirical}.  They are documented here for
completeness; every function has full \proglang{R} + \proglang{Python}
parity (the \proglang{Python} side is described in the morie
\proglang{Python} paper, this volume).

\paragraph{OTIS suite.}  Five callables operate on the Ontario
Tracking Information System public release.  \code{mrm_classify_mandela()}
(introduced in v0.1.14) classifies placements under the United Nations
Mandela Rules with three denominator conventions (row,
\code{individual_any}, \code{individual_cumulative}).
\code{mrm_otis_placement_concentration()} expands the b09 banded
per-individual placement counts via band midpoints and reports the
Hill-MLE Pareto exponent, the Gini coefficient, and the top-k\%
concentration share within each fiscal year and pooled.
\code{mrm_otis_seg_duration_km()} returns a Kaplan-Meier summary of
\texttt{NumberConsecutiveDays\_Segregation} with optional alert-flag
stratification and a fraction-above-Mandela-threshold column.
\code{mrm_otis_mortification_cooccurrence()} computes pairwise
Cram\'er's $V$ across the three b01 alert flags
(\texttt{MentalHealth\_Alert}, \texttt{SuicideRisk\_Alert},
\texttt{SuicideWatch\_Alert}).  \code{mrm_otis_region_locality()}
constructs the
\texttt{Region\_AtTimeOfPlacement}~$\times$~\texttt{Region\_MostRecentPlacement}
contingency table and reports the chi-square statistic, Cram\'er's
$V$, and the diagonal share (within-region staying).

\paragraph{TPS suite.}  Four callables operate on Toronto Police
Service open-data crime events.  \code{mrm_tps_levy_scaling()} fits
a Hill-MLE Pareto exponent to the haversine inter-incident step
lengths.  \code{mrm_tps_moran_clustering()} grids the lat/long
extent into a coarse raster, counts events per cell, and reports
the global Moran's $I$ z-score together with a DBSCAN summary at
$\varepsilon = 0.3$~km.
\code{mrm_tps_neighbourhood_recurrence_km()} returns a per-
\texttt{HOOD\_158} Kaplan-Meier summary of inter-event gap days.
\code{mrm_tps_kulldorff_scan()} implements Kulldorff's 1997
Poisson log-likelihood-ratio space-time scan
\citep{kulldorff1997} with a Monte-Carlo permutation test for
significance over $r \in \{1, 2, 3, 5, 8\}$~km candidate radii.

\paragraph{SIU suite.}  Three callables operate on Ontario Special
Investigations Unit case-level data.
\code{mrm_siu_case_to_decision_km()} reports a Kaplan-Meier summary
of the gap from \texttt{date\_of\_incident\_iso} to
\texttt{date\_of\_director\_decision\_iso} per
\texttt{police\_service}, with right-censoring of open cases.
\code{mrm_siu_per_service_rate()} cross-tabulates cases by service
and year.  \code{mrm_siu_outcome_classifier()} tabulates the
Director's decision categories per service.  Unlike OTIS (which
re-assigns its \texttt{UniqueIndividual\_ID} each fiscal year via
the \texttt{YYYY-XXXXX-SG} format and therefore admits no
cross-year tracking) and TPS (which exposes no offender
identifier by design), SIU provides per-case dates with a stable
jurisdiction identifier and so admits a valid time-to-outcome
analysis; the canonical example is the gap from
\texttt{date\_of\_incident\_iso} to
\texttt{date\_of\_director\_decision\_iso} per
\texttt{police\_service}, which our verification gives
as $120$ days (pooled, $n = 1{,}711$).

\paragraph{Longitudinal simulator.}  One callable,
\code{morie_simulate_longitudinal_panel()}, ties a synchronised
RNG, a VAR$(L)$ coefficient array with stationarity-preserving
spectral constraints, and structured-covariance multivariate
normal draws into a tidy long-format panel.  This is a clean-room
implementation of standard methods from
\citet{hamilton1994timeseries} and
\citet{diggle1994longitudinal}.

\paragraph{Dataset access paths.}  The empirical-workflow
callables are paired with three on-demand dataset fetchers so the
callables work out of the box without manual data wrangling: OTIS
via the Ontario Data Catalogue CKAN endpoints (resource IDs are
baked into \code{morie_dataset_catalog()} for a01, b01, b09 and
c11); TPS via an ArcGIS REST fetcher
(\code{morie_fetch_tps()}, nine public categories); and SIU via
\code{morie_fetch_siu()}, which scrapes the public Director's
Reports site on demand.  Four small reference samples ship in
\texttt{inst/extdata/} so the worked examples and the testthat
suite run offline.

\paragraph{Animated demo.}  A user can confirm the package is actually doing work, in the
same way DoubleML's console output does, by installing the R sibling and calling
\code{run\_treatment\_effects\_analysis()} with \code{verbose = TRUE}.
Per-step progress prints to the console via the \pkg{cli} package; the Python sibling provides an
equivalent animated demo via \code{python -m morie.demo}.

\section[Textbook-derived callables in morie.fn]{Textbook-derived callables in \texttt{morie.fn}}\label{sec:textbook-callables}

Beyond the causal-inference core and the empirical-workflow
callables of Section~\ref{sec:mrm-empirical-callables},
\pkg{morie} ships a curated catalogue of 275 textbook-derived
callables, organised into fourteen thematic suites under the
\code{morie.fn} namespace.  Each suite implements a coherent
slice of a canonical reference text in statistics, machine
learning, or signal processing; together they form a
pedagogical-and-applied corpus that complements the wrapper
functions documented above.  The design philosophy is uniform:
where a battle-tested library exists on either side of the
\proglang{R}/\proglang{Python} divide, the suite calls into it
(\code{spdep}, \code{gstat}, \code{survival}, \code{forecast},
\code{rugarch}, \code{signal}, \code{kernlab} on the
\proglang{R} side; \code{scipy}, \code{statsmodels},
\code{scikit-learn}, \code{pywavelets} on the \proglang{Python}
side); where the operation is light enough that a library would
introduce more weight than logic, the suite uses base \proglang{R}
matrix and vector primitives (or numpy on the Python side)
directly.  All 275 callables were verified for \proglang{R} +
\proglang{Python} parity at the most recent release (see the package
\code{NEWS} file).

The fourteen suites are enumerated in
Table~\ref{tab:textbook-suites}, with the short module mnemonic
that fixes the namespace, the number of callables shipped, the
canonical reference textbook, and one representative
\proglang{R} wrapper as illustration.

\begin{table}[t]
\centering
\footnotesize
\begin{tabular}{lr>{\raggedright\arraybackslash}p{4.5cm}>{\raggedright\arraybackslash}p{5.8cm}}
\toprule
Suite & $n$ & Reference & Representative wrapper \\
\midrule
\code{schabenberger} & 20 & \citet{SchabenbergerGotway2005} & \code{spatial\_autocorrelation()} \\
\code{kosorok} & 20 & \citet{Kosorok2008} & \code{kosorok\_empirical\_process()} \\
\code{ml\_foundations} & 20 & \citet{MohriRostamizadehTalwalkar2018,HastieTibshiraniFriedman2009} & \code{linear\_regression\_ols()} \\
\code{ghosal} & 20 & \citet{GhosalVanDerVaart2017} & \code{ghosal\_dirichlet\_posterior()} \\
\code{horowitz} & 20 & \citet{Horowitz2009} & \code{horowitz\_kernel\_regression()} \\
\code{montesinos} & 20 & \citet{MontesinosLopez2022} & \code{gblup\_full()} \\
\code{armstrong} & 20 & \citet{Armstrong2014SpatialVoting} & \code{ideal\_point\_recovery()} \\
\code{deep\_learning} & 20 & \citet{GoodfellowBengioCourville2016} & \code{forward\_pass\_dense()} \\
\code{missing\_stats} & 20 & QMC / copulas / EVT & \code{bootstrap\_ci()} \\
\code{llm\_arch} & 20 & \citet{Vaswani2017Attention,Dao2022FlashAttention,Su2021RoPE,ZhangSennrich2019RMSNorm} & \code{bpe\_tokenizer()} \\
\code{fauzi} & 15 & \citet{Fauzi2018Kernel} & \code{fauzi\_kdfe\_properties()} \\
\code{gibbons\_chakraborti} & 20 & \citet{GibbonsChakraborti2014} & \code{tolerance\_limits()} \\
\code{rangayyan} & 20 & \citet{Rangayyan2015} & \code{rangayyan\_fir\_filter()} \\
\code{time\_series\_advanced} & 20 & \citet{HyndmanAthanasopoulos2018,Tsay2010FinancialTimeSeries} & \code{garch\_fit()} \\
\bottomrule
\end{tabular}
\caption{The fourteen textbook-derived suites in \code{morie.fn}.
Total: 275 callables, all with \proglang{R}+\proglang{Python}
parity verified at v0.3.0.  The \code{missing\_stats} label
is a historical artefact: the spec covers resampling,
quasi-Monte-Carlo, copula, and extreme-value-theory material,
not missing-data methods proper; a true missing-data spec is
deferred to v0.3.x.}
\label{tab:textbook-suites}
\end{table}

\subsection{Naming convention}\label{sec:textbook-naming}

Each suite occupies a four-to-seven-character mnemonic
namespace.  On the \proglang{R} side every callable is exported
under \emph{both} its compact mnemonic and its full descriptive
name: in the \code{schabenberger} suite, for instance, both
\code{sptau()} and \code{spatial\_autocorrelation()} resolve to
the same function, because the full-name binding is created by
\code{spatial\_autocorrelation <- sptau}.  The mnemonic matches
the spec \code{name} field of the originating formula-index JSON
and the descriptive binding matches the \code{full\_name} field,
so a user can locate the originating specification and the
cross-language sibling without guessing.  The \proglang{Python} sibling uses the longer form
\code{morie.fn.<short>.<full\_name>(\ldots)} as documented in
\citet{Ruhela2026MoriePy}.

\subsection{Caveats and v0.3.x deferrals}\label{sec:textbook-caveats}

Numerical agreement between the \proglang{R} and \proglang{Python}
sides is bit-identical for deterministic callables: the
closed-form OLS and PCA in \code{ml\_foundations}, Moran's $I$
and ordinary kriging in \code{schabenberger}, the analytic
attention and softmax operators in \code{llm\_arch}, and similar.  For stochastic callables
(bootstrap and jackknife in \code{missing\_stats}, random-forest
and gradient-boosting ensembles, MCMC posteriors in
\code{ghosal}, the Bayesian-inference paths in
\code{time\_series\_advanced}), agreement is within Monte-Carlo
tolerance under matched random-number seeds.  The
\proglang{R} smoke suite confirms that all 319 \proglang{R}
files in the source tree source cleanly under base \proglang{R}
with only the declared \code{Suggests:} optional packages
installed; the \proglang{Python} smoke runs the \code{\_\_main\_\_}
canonical-test block in every shipped module.  Three honest
deferrals carry over from the v0.3.x tracker: an opt-in
deterministic mode for the stochastic callables, a true
missing-data suite to replace the misnamed \code{missing\_stats}
spec, and a tighter cross-language MLE-tolerance pin for the
GARCH/EGARCH and Bayesian fits where the two language
ecosystems' optimiser defaults differ.

\section{Disparity auditing}
\label{sec:fairness}

Release~0.8.0 adds \pkg{morie}'s fairness and disparity-audit surface:
a family of callables for \emph{auditing} risk-assessment and
predictive-policing systems for group disparity --- the package
measures whether such a system encodes disparate treatment; it does
not build one.  Six group-fairness metrics are exposed:
\code{fairness\_disparate\_impact} (the EEOC four-fifths rule),
\code{fairness\_demographic\_parity}, \code{fairness\_equalized\_odds},
\code{fairness\_average\_odds\_difference}, \code{fairness\_gini}, and
the composite \code{fairness\_bias\_amplification} after
\citet{Barman2026PredictivePolicingBias}.  The
\code{predpol\_calibration\_audit} callable generalises the Strategic
Subjects List district analysis of \citet{Lacherade2021PredPolChicago}:
it ranks areas by predicted risk against their realised outcome rate
and tests whether the rank disagreement tracks demographic
composition, while \code{predpol\_temporal\_audit} extends the audit
across time and cities.  Each callable returns the same structured
result object as the rest of the package, and all carry runnable
examples.  A model-agnostic explainability layer --- permutation
importance, partial dependence, accumulated local effects, ceteris
paribus, and Shapley attributions --- reimplements the COMPAS audit of
\citet{Biecek2021XAIStories}.  Every method is an independent
clean-room reimplementation from the cited works' published
descriptions; no third-party source code is incorporated.

\section{Conclusion}
\label{sec:conclusion}

The \proglang{R} package \pkg{morie} provides a curated
multi-domain toolkit for observational inference, with explicit
support for the MRM framework as a sociolegal application target.
The wrapper architecture around \pkg{DoubleML} offers a thin
function-style API; the eight thematic groups provide a unified
namespace for the broader observational-inference, signal-processing,
spatial, statistical-physics, and psychometric routines that
complement the causal core. The MRM modules implement the
framework of \citet{Ruhela2026MRM} on the five Canadian sources
through a common ingestion pipeline.

The \proglang{R} and \proglang{Python} packages are both released
under \code{AGPL-3.0-or-later}, a strong copyleft licence: any
modified version that is distributed, or offered to users over a
network, must publish its source.  The optional Linux kernel module
(\texttt{kernel-module/morie.c}) ships under GPL-2.0-only (the
kernel ABI requirement) and is not part of the package; papers and
documentation are \code{CC BY-NC-SA 4.0}.  The full per-component breakdown is in
\code{LICENSING.md} at the source repository
\url{https://github.com/hadesllm/morie}.  Contributions follow
the process documented in \code{CONTRIBUTING.md}.

\section*{Acknowledgements}

The author thanks Glenn McNamara for distribution-theory
mentorship that grounds the per-individual estimator ensemble;
Professor Angela Zorro Medina for the methodological lineage of
the framework's identification strategy; Jeroen Ooms for
maintaining the r-universe infrastructure on which the \pkg{morie}
nightly binaries are built; and the DoubleML team (Philipp Bach,
Malte S.\ Kurz, Victor Chernozhukov, Martin Spindler, Sven
Klaassen) for the BSD-3-Clause double-machine-learning library on
which the package's IRM and PLR wrappers are layered.

\paragraph*{AI assistance.}
This work was developed with substantial assistance from frontier
AI assistants.  The author retains full responsibility for the
code, the methods, and the scientific claims; AI assistance
accelerated implementation but does not change the attribution of
the work.  Anthropic's Claude family -- Opus, Sonnet, and Haiku
across the 4.x generation \citep{Anthropic2024ClaudeFamily} -- was
used extensively for code generation, refactoring, documentation,
code review, methodology discussion, and design iteration.  The
primary development interface was Claude Code
\citep{Anthropic2024ClaudeCode}, Anthropic's official
command-line agent.  The alignment foundations of Claude are
documented in \citet{Bai2022ConstitutionalAI}.  Use was supported
by Anthropic research-credit programs.  Google DeepMind's
Gemini 2.5 (Pro and Flash) models
\citep{GoogleDeepMind2024Gemini} on the Vertex AI platform
\citep{GoogleCloud2024VertexAI} were used for additional code
generation, cross-checking Claude-generated code, multi-modal
data analysis, and prototype evaluation.  Use was supported by
Google research-credit programs.  Symbolic code
navigation and structural editing across the codebase were
performed via the Serena MCP server \citep{Oraios2024Serena},
an open-source LSP-aggregator developed by Oraios AI that exposes
find-symbol, replace-symbol-body, and find-referencing-symbols
primitives to Claude Code over the Model Context Protocol.
\bibliography{refs}

\end{document}
